\begin{abstract}
Transmission spectroscopy uses the wavelength dependence of the transit depth
to probe the composition of exoplanet atmospheres. The spectrum is inferred
through a chain of detector calibration, background treatment, spectral
extraction and light-curve fitting. Independent pipelines can produce
different spectra from the same observation. Agreement between them cannot
establish accuracy, because the true spectrum of a real planet is unknown.
This project asks how accurately JWST/NIRISS SOSS transmission spectra can be
recovered from detector-level data, and whether fitting the transit and the
background jointly in the detector pixels reduces the biases of
extraction-first analyses. I have developed NOVA, a forward model that fits
the transit and the background directly in the detector pixels. It
propagates one transmission spectrum, shared by both spectral
orders, through an empirical instrument response, and fits it together with
the stellar continuum. A first, preliminary NOVA spectrum
of WASP-17~b shows features similar to those of the published spectrum, but
lies deeper at all wavelengths and most of all in the red. Its full
uncertainty is not yet validated. Because real data cannot settle such differences, I
built injection benchmarks on real SOSS data with a known transit. These
showed that, for one tested atmosphere, the ranking of two pipelines
reverses with the assumed fraction
of the observed light that belongs to the star, a quantity that the
WASP-17~b data do not determine. I am therefore building a benchmark from programme 4476, which observed a
star at its usual position and displaced from the detector strip. A cleaned
image of its starlight, with a known light curve and photon noise, will be
added to the raw data of the displaced exposure, and recovery tests will be
conditional on this construction. The rest of
the PhD will use this benchmark to score NOVA against published pipelines,
apply NOVA to five further SOSS visits and to repeated visits of active
stars, and extend it to NIRSpec and to secondary eclipses.
\end{abstract}
